% Reviewed translation: Section 3, PDF pages 311--314 / printed pages 297--300.
% The closing Remark 2.3 on PDF 311 belongs to the preceding file.
\MathBookSourcePage{311}
\section{非奇异解分支的逼近}
\label{sec:navier-nonsingular-branches}
\setcounter{equation}{0}

上一节已经证明, 一般而言 Navier--Stokes 方程有不止一个解, 除非数据
(即黏度和外力)满足非常严格的要求. 然而还可以证明, 在许多实际例子中,
这些解大多是孤立的, 即每个解都在某个邻域中唯一. 此外, 可以证明解连续依赖于黏度.
所以, 当黏度沿一个区间变化时, Navier--Stokes 方程的每个解描绘出一个孤立分支.
这尤其意味着分岔现象并不常见. 实践中经常遇到的这一情形, 在数学上由非奇异解分支的概念表达.

本节提出并分析一大类非线性问题(包括 Navier--Stokes 问题)的非奇异解分支的若干逼近.
这一分析以隐函数 Theorem 的一般形式为基础, 是 Brezzi, Rappaz 与
Raviart~\cite{BrezziRappazRaviart1980} 所发展更广理论的一种变体.
这里给出的版本由 Crouzeix~\cite{CrouzeixForthcoming} 建立.

\subsection{一个抽象框架}
\label{subsec:navier-nonsingular-abstract}

设 $X$ 和 $\mathscr X$ 是两个 Banach 空间, $\Lambda$ 是实直线
$\mathbb R$ 的一个紧区间. 给定一个 $C^p$ 映射($p\geq1$)
\[
F:(\lambda,u)\in\Lambda\times X
\longmapsto F(\lambda,u)\in\mathscr X,
\]
并希望求解方程
\begin{equation}\label{eq:navier-branch-abstract-problem}
F(\lambda,u)=0,
\end{equation}
即寻找满足 \eqref{eq:navier-branch-abstract-problem} 的
$(\lambda,u)\in\Lambda\times X$.

设 $\{(\lambda,u(\lambda));\lambda\in\Lambda\}$ 是方程
\eqref{eq:navier-branch-abstract-problem} 的一个解分支. 这意味着
\begin{equation}\label{eq:navier-branch-continuity}
\lambda\longmapsto u(\lambda)
\text{ is a continuous function from }\Lambda\text{ into }X,
\end{equation}
以及
\begin{equation}\label{eq:navier-branch-solution-identity}
F(\lambda,u(\lambda))=0.
\end{equation}
还假设这些解在如下意义下非奇异:
\begin{equation}\label{eq:navier-branch-nonsingularity}
D_uF(\lambda,u(\lambda))
\text{ is an isomorphism from }X\text{ onto }\mathscr X
\quad\forall\lambda\in\Lambda.
\end{equation}
由隐函数 Theorem, \eqref{eq:navier-branch-nonsingularity} 立即蕴含
$\lambda\mapsto u(\lambda)$ 是从 $\Lambda$ 到 $X$ 的一个 $C^p$ 函数.

\MathBookSourcePage{312}
作为基本例子, 说明速度--压力表述中的 Navier--Stokes 方程 Dirichlet 问题
\eqref{eq:navier-steady-nonhomogeneous-variational} 如何纳入上述框架. 首先置
\begin{equation}\label{eq:navier-branch-ns-spaces}
X=\mathscr X=H^1(\Omega)^N\times L_0^2(\Omega),
\end{equation}
并引入中间空间
\begin{equation}\label{eq:navier-branch-ns-data-space}
Y=H^{-1}(\Omega)^N\times
\left\{\mathbf g\in H^{1/2}(\Gamma)^N;
\ \int_{\Gamma_i}\mathbf g\cdot\mathbf n\,ds=0,
\ 0\leq i\leq p\right\}.
\end{equation}
定义线性算子 $T$ 如下: 给定 $(\mathbf f_*,\mathbf g_*)\in Y$, 以
\[
(\mathbf u_*,p_*)=T(\mathbf f_*,\mathbf g_*)\in X
\]
表示下列 Stokes 方程 Dirichlet 问题的解:
\begin{equation}\label{eq:navier-branch-stokes-solution-operator}
\left\{
\begin{aligned}
-\Delta\mathbf u_*+\operatorname{grad}p_*&=\mathbf f_*&&\text{in }\Omega,\\
\operatorname{div}\mathbf u_*&=0&&\text{in }\Omega,\\
\mathbf u_*|_\Gamma&=\mathbf g_*.
\end{aligned}
\right.
\end{equation}
最后, 对数据 $(\mathbf f,\mathbf g)\in Y$, 定义从
$\mathbb R_+\times X$ 到 $Y$ 的 $C^\infty$ 映射
\begin{equation}\label{eq:navier-branch-ns-nonlinear-map}
G:(\lambda,\mathbf v=(\mathbf v,q))
\longmapsto
G(\lambda,\mathbf v)
=\left(
\lambda\left(\sum_{j=1}^{N}v_j
\frac{\partial\mathbf v}{\partial x_j}-\mathbf f\right),
-\mathbf g
\right),
\end{equation}
并置
\begin{equation}\label{eq:navier-branch-ns-compound-map}
F(\lambda,\mathbf v)=\mathbf v+TG(\lambda,\mathbf v).
\end{equation}

\setcounter{statement}{0}
\begin{Lemma}\label{lem:navier-branch-ns-equivalence}
一对 $(\mathbf u,p)\in H^1(\Omega)^N\times L_0^2(\Omega)$ 是问题
\eqref{eq:navier-steady-system},
\eqref{eq:navier-steady-nonhomogeneous-dirichlet} 的解, 当且仅当
$\boldsymbol u=(\mathbf u,p/\nu)$ 是
\eqref{eq:navier-branch-abstract-problem} 的一个解, 其中
$\lambda=1/\nu$, 空间 $X$ 和 $\mathscr X$ 由
\eqref{eq:navier-branch-ns-spaces} 定义, 复合映射 $F$ 由
\eqref{eq:navier-branch-ns-compound-map} 定义.
\end{Lemma}
\begin{Proof}
方程 \eqref{eq:navier-steady-system} 可等价地写成
\[
\begin{aligned}
-\Delta\mathbf u+\operatorname{grad}(p/\nu)
&=\frac1\nu\left(\mathbf f-
\sum_{j=1}^{N}u_j\frac{\partial\mathbf u}{\partial x_j}\right)
&&\text{in }\Omega,\\
\operatorname{div}\mathbf u&=0&&\text{in }\Omega,\\
\mathbf u&=\mathbf g&&\text{on }\Gamma.
\end{aligned}
\]
或由 \eqref{eq:navier-branch-stokes-solution-operator} 写成
\[
(\mathbf u,p/\nu)
=T\left[
\frac1\nu\left(\mathbf f-
\sum_{j=1}^{N}u_j\frac{\partial\mathbf u}{\partial x_j}\right),
\mathbf g\right].
\]
结论于是由 \eqref{eq:navier-branch-ns-nonlinear-map} 和
\eqref{eq:navier-branch-ns-compound-map} 得到.
\end{Proof}

\MathBookSourcePage{313}
接着注意, 对所有 $\boldsymbol v=(\mathbf v,q)\in X$,
\[
D_{\boldsymbol u}G(\lambda,\boldsymbol u)\cdot\boldsymbol v
=\left(
\lambda\sum_{j=1}^{N}\left(
u_j\frac{\partial\mathbf v}{\partial x_j}
+v_j\frac{\partial\mathbf u}{\partial x_j}\right),0
\right).
\]
因此在本例中, $\boldsymbol u=(\mathbf u,p/\nu)$ 是
\eqref{eq:navier-branch-abstract-problem} 的一个非奇异解, 或等价地
$D_{\boldsymbol u}F(1/\nu,\boldsymbol u)$ 是 $X$ 的一个同构, 当且仅当对每个
$\boldsymbol w=(\mathbf w,\sigma)\in X$, 存在唯一的
$\boldsymbol v=(\mathbf v,q)\in X$ 满足
\[
\begin{aligned}
-\Delta\mathbf v+\operatorname{grad}q
+\frac1\nu\sum_{j=1}^{N}\left(
u_j\frac{\partial\mathbf v}{\partial x_j}
+v_j\frac{\partial\mathbf u}{\partial x_j}\right)
&=-\Delta\mathbf w+\operatorname{grad}\sigma&&\text{in }\Omega,\\
\operatorname{div}\mathbf v&=\operatorname{div}\mathbf w&&\text{in }\Omega,\\
\mathbf v&=\mathbf w&&\text{on }\Gamma.
\end{aligned}
\]
这一问题可以化简. 一方面, Cattabriga~\cite{Cattabriga1961} 证明的
Theorem~I.5.4 的直接推广保证: 对每个
$\mathbf f\in H^{-1}(\Omega)^N$, $\mu\in L^2(\Omega)$ 和
$\mathbf g\in H^{1/2}(\Gamma)^N$ 满足相容条件
\[
\int_\Omega\mu\,dx=\int_\Gamma\mathbf g\cdot\mathbf n\,ds,
\]
存在唯一的 $\boldsymbol w=(\mathbf w,\sigma)\in X$ 满足
\[
\begin{aligned}
-\Delta\mathbf w+\operatorname{grad}\sigma&=\mathbf f&&\text{in }\Omega,\\
\operatorname{div}\mathbf w&=\mu&&\text{in }\Omega,\\
\mathbf w&=\mathbf g&&\text{on }\Gamma.
\end{aligned}
\]
另一方面, 当 $\mu$ 和 $\mathbf g$ 满足上述条件时,
Lemma~\ref{lem:sec2-div-free-lifting} 的一个推广表明存在
$\mathbf v_0\in H^1(\Omega)^N$, 使
\[
\operatorname{div}\mathbf v_0=\mu,
\qquad
\mathbf v_0|_\Gamma=\mathbf g.
\]
因此利用这两个结果可以证明, $\boldsymbol u=(\mathbf u,p/\nu)$ 是
\eqref{eq:navier-branch-abstract-problem} 的一个非奇异解, 当且仅当齐次线性化问题
\begin{equation}\label{eq:navier-branch-linearized-problem}
\left\{
\begin{aligned}
-\nu\Delta\mathbf v
+\sum_{j=1}^{N}\left(
u_j\frac{\partial\mathbf v}{\partial x_j}
+v_j\frac{\partial\mathbf u}{\partial x_j}\right)
+\operatorname{grad}q&=\mathbf f_*&&\text{in }\Omega,\\
\operatorname{div}\mathbf v&=0&&\text{in }\Omega,\\
\mathbf v&=0&&\text{on }\Gamma
\end{aligned}
\right.
\end{equation}
对每个 $\mathbf f_*\in H^{-1}(\Omega)^N$ 有唯一解
$\boldsymbol v=(\mathbf v,q)\in X$, 且映射
$\mathbf f_*\mapsto\boldsymbol v$ 从 $H^{-1}(\Omega)^N$ 到 $X$ 连续.

尽管这里不准备给出精确陈述, 仍可证明 Navier--Stokes 方程 Dirichlet 问题的解
``一般而言''是非奇异的. 这里只推出下面的简单结果.

\MathBookSourcePage{314}
\setcounter{statement}{1}
\begin{Lemma}\label{lem:navier-branch-small-data-nonsingular}
假设 Theorem~\ref{thm:navier-steady-nonhomogeneous-uniqueness} 的条件成立.
若 $(\mathbf u,p)$ 是问题
\eqref{eq:navier-steady-nonhomogeneous-variational} 的唯一解, 则
$\boldsymbol u=(\mathbf u,p/\nu)$ 是
\eqref{eq:navier-branch-abstract-problem} 的一个非奇异解.
\end{Lemma}
\begin{Proof}
先利用 Subsection~\ref{subsec:navier-steady-velocity-pressure} 的记号.
问题 \eqref{eq:navier-branch-linearized-problem} 的一个变分形式是求
$(\mathbf v,q)\in H^1(\Omega)^N\times L_0^2(\Omega)$, 使
\begin{equation}\label{eq:navier-branch-linearized-variational}
\left\{
\begin{aligned}
a(\mathbf u;\mathbf v,\mathbf z)
+a_1(\mathbf v;\mathbf u,\mathbf z)
-(q,\operatorname{div}\mathbf z)
&=\langle\mathbf f_*,\mathbf z\rangle
&&\forall\mathbf z\in H_0^1(\Omega)^N,\\
\operatorname{div}\mathbf v&=0&&\text{in }\Omega,\\
\mathbf v&=\mathbf g_*&&\text{on }\Gamma.
\end{aligned}
\right.
\end{equation}
要验证 \eqref{eq:navier-branch-linearized-variational} 有唯一解, 只需证明双线性形式
\[
(\mathbf v,\mathbf z)\longmapsto
c(\mathbf v,\mathbf z)
=a(\mathbf u;\mathbf v,\mathbf z)
+a_1(\mathbf v;\mathbf u,\mathbf z)
\]
是 $V$-椭圆的. 由 \eqref{eq:navier-steady-convection-cancellation},
\[
c(\mathbf v,\mathbf v)
=\nu|\mathbf v|_{1,\Omega}^2
+a_1(\mathbf v;\mathbf u,\mathbf v)
\qquad\forall\mathbf v\in V.
\]
设 $\nu>\nu_0$, 其中 $\nu_0$ 由
\eqref{eq:navier-steady-critical-viscosity} 定义. 则存在
$\mathbf u_0\in H^1(\Omega)^N$, 使
\[
\operatorname{div}\mathbf u_0=0,
\qquad \mathbf u_0|_\Gamma=\mathbf g,
\qquad
\nu>\rho(\mathbf u_0)
+\bigl(\mathcal N\|l(\mathbf f;\mathbf u_0)\|_{V'}\bigr)^{1/2}.
\]
置 $\mathbf u=\mathbf u_0+\mathbf w$. 由
\eqref{eq:navier-steady-convection-constant} 和
\eqref{eq:navier-steady-lifting-rho},
\[
|a_1(\mathbf v;\mathbf u,\mathbf v)|
\leq|a_1(\mathbf v;\mathbf u_0,\mathbf v)|
+|a_1(\mathbf v;\mathbf w,\mathbf v)|
\leq\{\rho(\mathbf u_0)+\mathcal N|\mathbf w|_{1,\Omega}\}
|\mathbf v|_{1,\Omega}^2.
\]
由 Remark~\ref{rem:navier-steady-lifted-bound},
\[
|\mathbf w|_{1,\Omega}
\leq\frac{\|l(\mathbf f;\mathbf u_0)\|_{V'}}
{\nu-\rho(\mathbf u_0)}.
\]
所以
\[
\begin{aligned}
c(\mathbf v,\mathbf v)
&\geq\left(\nu-\rho(\mathbf u_0)
-\frac{\mathcal N\|l(\mathbf f;\mathbf u_0)\|_{V'}}
{\nu-\rho(\mathbf u_0)}\right)|\mathbf v|_{1,\Omega}^2\\
&=(\nu-\rho(\mathbf u_0))
\left(1-
\frac{\mathcal N\|l(\mathbf f;\mathbf u_0)\|_{V'}}
{(\nu-\rho(\mathbf u_0))^2}\right)|\mathbf v|_{1,\Omega}^2.
\end{aligned}
\]
因此椭圆性成立.
\end{Proof}

后续各小节将主要研究 Navier--Stokes 方程 Dirichlet 问题的非奇异解分支
\[
\{(\lambda,\boldsymbol u(\lambda)
=(\mathbf u(\lambda),\lambda p(\lambda)));
\ \lambda\in\Lambda\}
\]
的逼近, 其中参数 $\lambda=1/\nu$ 起 Reynolds 数的作用.
